\section{搜索与多候选选择}

\subsection{自一致采样与多候选}
在 search 层面，自一致性（Self-Consistency）成为核心武器：允许模型生成多条候选，再用多数票或 consistency score 挑出最稳定的 final answer。如图 \ref{fig:self-consistency} 所见，随着候选数量从 10 织到 40，accuracy 仍在提升，远甩基于 log-prob 的排序。

\begin{figure}[H]
\centering
\includegraphics[width=0.92\textwidth]{frame-self-consistency.jpg}
\caption{Self-Consistency 曲线对比 log-prob 排序与其他 sampling 策略。\protect\footnotemark}
\label{fig:self-consistency}
\end{figure}
\footnotetext{视频画面时间区间：00:45:05--00:50:05，图中多样化采样与投票显著优于单一 log-prob 排序。}

\begin{knowledgebox}{多候选采样三要素}
1. 多样性：高 temperature、nucleus sampling 与 prompt 顺序扰动保证 candidate 之间差异；2. 选择策略：多数票优于 log-prob，除非为某任务训练了强 verifier；3. 资源分配：multi-sample 仅在需要更高准确率、允许延迟的场景打开。
\end{knowledgebox}

\subsection{Self-Consistency 与 Search 的协同}
Self-Consistency 为 search 提供了更宽的 width，而 beam search/Tree-of-Thought 通过 stepwise scoring 控制 depth。Chen 在 47 分 02 秒指出，理想的流程是先用 sampling 生成 20~40 个 candidate，再用 search + verifier 判断哪条路径可继续深化，最后用 consistency 投票稳定最终 answer。简单使用 majority vote 虽然可靠，但要在 inference-time 控制成本，仍需和 search 彼此配合，一边确保候选多样，一边用 verifier 在 early stage 就截断 bad paths。

\begin{warningbox}{不要只靠其中一个模块}
如果只运行 Self-Consistency 而不加任何 search，可能会在 token 成本上耗尽却继续重复相似错误；反之，如果只跑 search 而不加 sampling，多数时候只能走到局部最优。因此必须在 inference-time pipeline 中把 sample 生成、search 拓展、consistency 投票三者串联。
\end{warningbox}

\subsection{验证器与排序基线}
若训练出 stepwise verifier，可在生成早期筛掉低质分支——DeepMind 的 AlphaCode trace 就是一个例子，它自述 code execution 并在每一步给分，进而影响后续 search。标量 verifier 要兼顾稳定性与可解释性，否则可能把复杂正确答案误判为“不可解释的黑盒”。

\begin{warningbox}{弱 verifier 的陷阱}
若 verifier 过度偏向短输出或训练集模板，它会把多步骤推理判定为“低分”而舍弃，反而让 self-consistency 收敛到一组常见但错误的解。要让 verifier 同时敏感 accuracy 与 reasoning depth。
\end{warningbox}

\subsection{搜索的工程线索}
在 52 分 50 秒的幻灯片中，Chen 进一步把 sampling、search、verifier 的输出量化：采样 count、candidate score、voter margin 等指标要在 inference-time logs 中持续扫描。一个常见 pipeline 是用 high-temperature sampling 生成 30 条 response，用 stepwise verifier 给每条 response 赋分，再用 majority vote 冠军最后解。每个环节都应有对应的监控表格，以便在真实部署中抓到“今天采样 diversity 缺失”或“verifier drift”的 warning。

\begin{knowledgebox}{搜索流水线指标}
1. Sampling depth：候选数量与 diversity；2. Verifier quality：stepwise score 趋势与稳定性；3. Voting stability：majority margin 或 consistency rate，用于判断是否需要额外 search。
\end{knowledgebox}

\subsection{本章小结}
自一致性配合多样化采样可以让模型从 mistakes 中恢复，但要结合高质量 verifier 排序。整套机制需要在 accuracy、token budget 与 latency 之间做权衡。

